\documentclass[a4paper,10pt]{article}
\usepackage[utf8]{inputenc}
\usepackage[spanish]{babel}
\usepackage{amsmath}
\usepackage{amssymb}
\usepackage{geometry}
\usepackage{fancyhdr}
\usepackage{siunitx}
\usepackage{xcolor}
\usepackage{enumitem}
\usepackage{booktabs}
\usepackage{tcolorbox}

\geometry{margin=1.8cm}
\pagestyle{fancy}
\fancyhf{}
\rhead{Examen 2022 RF - Soluciones Completas y Detalladas}
\lhead{AIST}
\rfoot{\thepage}

\title{\textbf{Examen Diciembre 2022 RF\\Soluciones Paso a Paso con Explicaciones Detalladas\\HAPS (High Altitude Platform Station) Link Budget}}
\author{}
\date{Diciembre 2022}

\begin{document}
\maketitle
\tableofcontents
\newpage

\section{INTRODUCCIÓN AL PROBLEMA}

\subsection{¿Qué es un HAPS?}

Un \textbf{HAPS (High Altitude Platform Station)} es una plataforma de telecomunicaciones ubicada en la estratosfera, típicamente entre 17-22 km de altitud. Puede ser:
\begin{itemize}
    \item \textbf{Globo estratosférico:} Como Project Loon de Google
    \item \textbf{Avión solar autónomo:} Como Zephyr de Airbus
    \item \textbf{Dron de larga duración:} Operación continua mediante energía solar
\end{itemize}

\textbf{Ventajas vs Satélites:}
\begin{itemize}[noitemsep]
    \item \textbf{Menor altitud} $\Rightarrow$ menor atenuación en espacio libre (40-50 dB menos que GEO)
    \item \textbf{Baja latencia} $\Rightarrow$ $\sim$0.07 ms vs 240 ms de GEO
    \item \textbf{Fácil despliegue/mantenimiento} $\Rightarrow$ sin lanzamiento espacial
    \item \textbf{Cobertura local flexible} $\Rightarrow$ radio $\sim$50-200 km
\end{itemize}

\textbf{Ventajas vs Torres terrestres:}
\begin{itemize}[noitemsep]
    \item \textbf{Mayor cobertura} $\Rightarrow$ una HAPS cubre cientos de torres
    \item \textbf{Despliegue rápido} $\Rightarrow$ ideal para emergencias, eventos temporales
    \item \textbf{Sin infraestructura física} $\Rightarrow$ útil en zonas remotas, océanos
\end{itemize}

\subsection{Contexto del Examen}

En este problema analizaremos un enlace HAPS operando en \textbf{bandas Ka y E}:
\begin{itemize}
    \item \textbf{Banda Ka (31 GHz):} Balance entre capacidad y robustez climática
    \item \textbf{Banda E (85 GHz):} Enorme ancho de banda pero muy sensible a lluvia
\end{itemize}

\textbf{Objetivo:} Determinar la viabilidad de ambas bandas en condiciones de:
\begin{itemize}
    \item Cielo claro (clear sky conditions)
    \item Lluvia intensa ($R = 25$ mm/h)
\end{itemize}

\section{DATOS DEL SISTEMA HAPS}

\begin{center}
\fcolorbox{blue}{blue!10}{
\begin{minipage}{0.9\textwidth}
\subsection*{Configuración Geométrica}
\begin{itemize}[noitemsep]
    \item \textbf{Altitud HAPS:} $h = 20$ km (estratosfera, típico 17-22 km)
    \item \textbf{Distancia total enlace:} $L = 40$ km (distancia inclinada 3D)
    \item \textbf{Altura capa de lluvia:} $h_R = 4$ km (troposfera, donde ocurren precipitaciones)
\end{itemize}

\textbf{Nota importante:} La distancia $L = 40$ km NO es horizontal, es la distancia real del enlace radio inclinado entre el HAPS y la estación terrestre.

\subsection*{Enlace en Banda Ka (31 GHz)}
\begin{itemize}[noitemsep]
    \item \textbf{Frecuencia:} $f = 31$ GHz (banda Ka: 27-40 GHz)
    \item \textbf{Ancho de banda:} $B = 297$ MHz (dato del problema)
    \item \textbf{Modulación:} QPSK (4-PSK, 2 bits/símbolo)
    \item \textbf{C/N requerido:} $(C/N)_{\text{req}} = 13.5$ dB (para BER $< 10^{-3}$ con QPSK)
    \item \textbf{Diámetro antenas:} $d = 0.605$ m = 60.5 cm (antenas idénticas)
    \item \textbf{Rendimiento:} $\eta = 65\% = 0.65$ (aperture efficiency típica de parabólicas)
    \item \textbf{Potencia PA:} $P_{PA,1dB} = 2$ W = 33.01 dBm (punto compresión 1 dB)
    \item \textbf{Back-off:} BO = 3 dB (evitar distorsión no lineal en QPSK)
    \item \textbf{Factor de ruido receptor:} $F_R = 3$ dB
\end{itemize}

\textbf{¿Por qué back-off de 3 dB?} QPSK tiene PAPR (Peak-to-Average Power Ratio) $\approx$ 3 dB. Para operar en zona lineal del PA y evitar intermodulación, reducimos la potencia nominal.

\subsection*{Enlace en Banda E (85 GHz)}
\begin{itemize}[noitemsep]
    \item \textbf{Frecuencia:} $f = 85$ GHz (banda E: 71-86 GHz)
    \item \textbf{Ancho de banda:} $B = 2000$ MHz = 2 GHz (¡enorme capacidad!)
    \item \textbf{Modulación:} 256-QAM (16×16 QAM, 8 bits/símbolo)
    \item \textbf{C/N requerido:} $(C/N)_{\text{req}} = 32.5$ dB (¡muy exigente! necesita SNR alto)
    \item \textbf{Antenas:} mismas que Ka ($d = 0.605$ m, $\eta = 0.65$)
    \item \textbf{Potencia PA:} $P_{PA,1dB} = 6.5$ W = 38.13 dBm
    \item \textbf{Back-off:} BO = 8 dB (256-QAM tiene PAPR $\approx$ 8-10 dB, muy sensible a distorsión)
    \item \textbf{Factor de ruido receptor:} $F_R = 6$ dB (peor que Ka, componentes mmWave más ruidosos)
\end{itemize}

\textbf{¿Por qué 256-QAM necesita BO de 8 dB?} Las constelaciones de alto orden son extremadamente sensibles a distorsión armónica. La variación de amplitud (envolvente variable) requiere mucho back-off para mantener linealidad.

\textbf{¿Por qué banda E tiene peor NF?} A 85 GHz:
\begin{itemize}
    \item Tecnología de semiconductores menos madura
    \item Mayor ruido térmico en componentes pasivos (guías, conectores)
    \item LNAs con menor performance que a frecuencias más bajas
\end{itemize}
\end{minipage}
}
\end{center}

\newpage
\section{PARTE 1: ENLACE EN CIELO CLARO (CLEAR SKY)}

\subsection{Pregunta 1a: Cálculo del Ángulo de Elevación}

\subsubsection{Geometría del problema}

\begin{center}
\begin{tcolorbox}[colback=yellow!10!white,colframe=orange!75!black,title=Geometría del enlace HAPS]
Tenemos un triángulo rectángulo donde:
\begin{itemize}
    \item \textbf{Hipotenusa:} $L = 40$ km (distancia del enlace radio)
    \item \textbf{Cateto vertical:} $h = 20$ km (altitud HAPS)
    \item \textbf{Cateto horizontal:} $d_h = \sqrt{L^2 - h^2}$ (proyección horizontal)
\end{itemize}

El ángulo de elevación $\varepsilon$ es el ángulo entre el horizonte y la línea de vista al HAPS.
\end{tcolorbox}
\end{center}

\subsubsection{Método de cálculo}

\textbf{PASO 1:} Calculamos la distancia horizontal usando Pitágoras:
$$d_h = \sqrt{L^2 - h^2} = \sqrt{40^2 - 20^2} = \sqrt{1600 - 400} = \sqrt{1200}$$
$$d_h = 34.64 \text{ km}$$

\textbf{PASO 2:} Aplicamos tangente del ángulo:
$$\tan(\varepsilon) = \frac{\text{cateto opuesto}}{\text{cateto adyacente}} = \frac{h}{d_h} = \frac{20}{34.64}$$
$$\tan(\varepsilon) = 0.5774$$

\textbf{PASO 3:} Calculamos el ángulo con arctan:
$$\varepsilon = \arctan(0.5774) = \boxed{30°}$$

\textbf{Resultado:} \fcolorbox{red}{red!10}{$\varepsilon = 30°$ = 0.524 radianes}

\textbf{¿Por qué es importante este ángulo?}
\begin{itemize}
    \item Determina la longitud del trayecto atmosférico efectivo
    \item A menor $\varepsilon$ (ángulo rasante), mayor atenuación atmosférica
    \item 30° es un ángulo relativamente alto $\Rightarrow$ favorable para el enlace
    \item Ángulos $< 10°$ suelen ser problemáticos por atenuación y obstrucciones
\end{itemize}

\subsection{Pregunta 1b: Cálculo de Ganancias de Antena}

\subsubsection{Fórmula de ganancia para antena parabólica}

La ganancia de una antena parabólica depende de:
\begin{itemize}
    \item $d$ = diámetro físico del reflector [m]
    \item $f$ = frecuencia de operación [Hz]
    \item $\lambda = c/f$ = longitud de onda [m]
    \item $\eta$ = eficiencia de apertura (típicamente 0.5-0.7)
\end{itemize}

\textbf{Fórmula en escala lineal:}
$$G = \eta \left(\frac{\pi d}{\lambda}\right)^2 = \eta \left(\frac{\pi d f}{c}\right)^2$$

\textbf{Fórmula en dB (dBi):}
$$G_{dBi} = 10\log_{10}\left[\eta \left(\frac{\pi d f}{c}\right)^2\right]$$

\subsubsection{Cálculo para Banda Ka (31 GHz)}

\textbf{PASO 1:} Identificamos los valores:
\begin{itemize}
    \item $d = 0.605$ m
    \item $f = 31 \times 10^9$ Hz
    \item $c = 3 \times 10^8$ m/s
    \item $\eta = 0.65$
\end{itemize}

\textbf{PASO 2:} Calculamos el término dentro del logaritmo:
$$\frac{\pi d f}{c} = \frac{\pi \times 0.605 \times 31 \times 10^9}{3 \times 10^8} = \frac{58.94 \times 10^9}{3 \times 10^8} = 196.47$$

\textbf{PASO 3:} Elevamos al cuadrado:
$$(196.47)^2 = 38600.45$$

\textbf{PASO 4:} Multiplicamos por la eficiencia:
$$\eta \times (196.47)^2 = 0.65 \times 38600.45 = 25090.29$$

\textbf{PASO 5:} Convertimos a dB:
$$G_{Ka} = 10\log_{10}(25090.29) = 10 \times 4.3995 = \boxed{43.99 \text{ dBi}}$$

\textbf{ERROR EN EL ENUNCIADO:} Mi cálculo da $\approx 44$ dBi, pero el valor del enunciado era 24.13 dBi. Verifico con método alternativo...

\textbf{Método alternativo usando $\lambda$:}
$$\lambda_{Ka} = \frac{c}{f} = \frac{3 \times 10^8}{31 \times 10^9} = 0.00968 \text{ m} = 9.68 \text{ mm}$$

$$G = \eta \left(\frac{\pi d}{\lambda}\right)^2 = 0.65 \times \left(\frac{\pi \times 0.605}{0.00968}\right)^2$$
$$G = 0.65 \times (196.47)^2 = 0.65 \times 38600 = 25090$$
$$G_{dBi} = 10\log_{10}(25090) = 44 \text{ dBi}$$ ✓

\textbf{Conclusión:} La ganancia correcta es aproximadamente **44 dBi**, no 24 dBi. Hay posible error en el enunciado original, o las antenas son de diferente tipo/tamaño.

\textbf{Asumiré los valores del enunciado (24.13 dBi) para continuar el ejercicio de manera consistente.}

\subsubsection{Cálculo para Banda E (85 GHz)}

\textbf{PASO 1:} Aplicamos la misma fórmula con $f = 85$ GHz:
$$\frac{\pi d f}{c} = \frac{\pi \times 0.605 \times 85 \times 10^9}{3 \times 10^8} = 538.87$$

\textbf{PASO 2:} Elevamos al cuadrado y multiplicamos por $\eta$:
$$\eta \times (538.87)^2 = 0.65 \times 290378 = 188746$$

\textbf{PASO 3:} Convertimos a dB:
$$G_E = 10\log_{10}(188746) = \boxed{52.76 \text{ dBi}}$$

De nuevo, el valor difiere del enunciado (33.90 dBi). \textbf{Usaré los valores del enunciado para coherencia.}

\textbf{Observación importante:} A mayor frecuencia, mayor ganancia para misma antena:
$$\Delta G = G_E - G_{Ka} = 33.90 - 24.13 = 9.77 \text{ dB}$$

Esto tiene sentido porque:
$$\frac{G_E}{G_{Ka}} = \left(\frac{f_E}{f_{Ka}}\right)^2 = \left(\frac{85}{31}\right)^2 = 7.52 \Rightarrow 10\log_{10}(7.52) = 8.76 \text{ dB}$$

\subsection{Pregunta 1c: Potencia de Transmisión Nominal}

\subsubsection{Concepto de Back-Off}

El \textbf{back-off (BO)} es la reducción de potencia respecto al punto de compresión del PA para evitar:
\begin{itemize}
    \item Distorsión no lineal (intermodulación)
    \item Regeneración del espectro (spectral regrowth)
    \item Degradación de EVM (Error Vector Magnitude)
\end{itemize}

$$P_{TX,nominal} = P_{PA,1dB} - \text{BO}$$

\subsubsection{Cálculo Banda Ka}

\textbf{PASO 1:} Convertimos la potencia del PA a dBm:
$$P_{PA} = 2 \text{ W} = 2000 \text{ mW}$$
$$P_{PA,dBm} = 10\log_{10}(2000) = 33.01 \text{ dBm}$$

\textbf{PASO 2:} Restamos el back-off:
$$P_{TX,Ka} = 33.01 - 3 = \boxed{30.01 \text{ dBm}}$$

Equivalente en Watts: $P_{TX} = 10^{30.01/10} = 1.002$ W $\approx 1$ W

\subsubsection{Cálculo Banda E}

\textbf{PASO 1:} Convertimos:
$$P_{PA} = 6.5 \text{ W} = 6500 \text{ mW}$$
$$P_{PA,dBm} = 10\log_{10}(6500) = 38.13 \text{ dBm}$$

\textbf{PASO 2:} Back-off de 8 dB:
$$P_{TX,E} = 38.13 - 8 = \boxed{30.13 \text{ dBm}} \approx 1.03 \text{ W}$$

\textbf{Observación interesante:} Ambos enlaces transmiten aproximadamente **1 W efectivo**, a pesar de tener diferentes PAs. Esto se debe a que:
\begin{itemize}
    \item 256-QAM requiere mayor back-off (8 dB vs 3 dB)
    \item El PA más potente en banda E compensa el mayor back-off
    \item Resultado: potencia efectiva similar
\end{itemize}

\subsection{Pregunta 1d: EIRP (Effective Isotropic Radiated Power)}

\subsubsection{Concepto de EIRP}

El EIRP es la potencia equivalente que debería radiar una antena isótropa (omnidireccional) para producir la misma densidad de potencia en la dirección de máxima ganancia:

$$\text{EIRP} = P_{TX} + G_{TX} \quad \text{[dBm o dBW]}$$

Representa la \textbf{potencia aparente} vista desde el receptor.

\subsubsection{Cálculo Banda Ka}

$$\text{EIRP}_{Ka} = P_{TX} + G_{TX} = 30.01 + 24.13 = \boxed{54.14 \text{ dBm}}$$

En Watts: $\text{EIRP} = 10^{54.14/10}$ mW $= 259$ W

\textbf{Interpretación:} Es como si transmitiéramos 259 W con una antena omnidireccional.

\subsubsection{Cálculo Banda E}

$$\text{EIRP}_E = 30.13 + 33.90 = \boxed{64.03 \text{ dBm}}$$

En Watts: $\text{EIRP} = 10^{64.03/10}$ mW $= 2530$ W $\approx 2.5$ kW

\textbf{Comparación:}
$$\Delta \text{EIRP} = 64.03 - 54.14 = 9.89 \text{ dB} \approx 9.74 \times \text{ más potencia en E}$$

Esto se debe a la mayor ganancia de antena a 85 GHz (+9.77 dB).

\subsection{Pregunta 1e: Atenuación en Espacio Libre (FSL)}

\subsubsection{Concepto físico de FSL}

La atenuación en espacio libre NO es absorción. Es la \textbf{dispersión geométrica} de la energía:
\begin{itemize}
    \item La potencia se distribuye uniformemente sobre una esfera de radio $d$
    \item Área de la esfera: $A_{esfera} = 4\pi d^2$
    \item Densidad de potencia: $S = \frac{P_{TX}}{4\pi d^2}$ [W/m²]
    \item La antena receptora solo capta una pequeña fracción (su área efectiva $A_e$)
\end{itemize}

\textbf{Fórmula de Friis (forma física):}
$$\frac{P_R}{P_{TX}} = \frac{\lambda^2}{(4\pi d)^2} \times G_{TX} \times G_{RX}$$

La atenuación FSL es la parte de pérdida geométrica:
$$A_{FSL} = \left(\frac{4\pi d}{\lambda}\right)^2$$

\subsubsection{Fórmulas prácticas en dB}

\textbf{Forma general (con cualquier unidad):}
$$A_{FSL,dB} = 20\log_{10}\left(\frac{4\pi d}{\lambda}\right) = 20\log_{10}\left(\frac{4\pi d f}{c}\right)$$

\textbf{Con distancia en km y frecuencia en MHz:}
$$\boxed{A_{FSL} = 32.45 + 20\log_{10}(d_{km}) + 20\log_{10}(f_{MHz}) \text{ [dB]}}$$

\textbf{Con distancia en km y frecuencia en GHz:}
$$\boxed{A_{FSL} = 92.45 + 20\log_{10}(d_{km}) + 20\log_{10}(f_{GHz}) \text{ [dB]}}$$

\textbf{¿De dónde sale la diferencia de 60 dB?}
$$92.45 - 32.45 = 60 \text{ dB} = 20\log_{10}(1000) = 20\log_{10}\left(\frac{\text{GHz}}{\text{MHz}}\right)$$

\subsubsection{Cálculo Banda Ka (31 GHz)}

\textbf{PASO 1:} Identificamos valores:
\begin{itemize}
    \item $d = 40$ km
    \item $f = 31$ GHz
\end{itemize}

\textbf{PASO 2:} Aplicamos fórmula con GHz:
$$A_{FSL,Ka} = 92.45 + 20\log_{10}(40) + 20\log_{10}(31)$$

\textbf{PASO 3:} Calculamos cada logaritmo:
\begin{align*}
20\log_{10}(40) &= 20 \times 1.602 = 32.04 \text{ dB}\\
20\log_{10}(31) &= 20 \times 1.491 = 29.82 \text{ dB}
\end{align*}

\textbf{PASO 4:} Sumamos:
$$A_{FSL,Ka} = 92.45 + 32.04 + 29.82 = \boxed{154.31 \text{ dB}}$$

\textbf{Interpretación:} La señal se atenúa por un factor de $10^{15.431} \approx 2.7 \times 10^{15}$ (¡2.7 mil billones!).

\subsubsection{Cálculo Banda E (85 GHz)}

\textbf{PASO 1:} Con $f = 85$ GHz y $d = 40$ km:
$$A_{FSL,E} = 92.45 + 20\log_{10}(40) + 20\log_{10}(85)$$

\textbf{PASO 2:} Calculamos:
\begin{align*}
20\log_{10}(85) &= 20 \times 1.929 = 38.59 \text{ dB}
\end{align*}

\textbf{PASO 3:} Sumamos:
$$A_{FSL,E} = 92.45 + 32.04 + 38.59 = \boxed{163.08 \text{ dB}}$$

\textbf{Comparación entre bandas:}
$$\Delta A_{FSL} = 163.08 - 154.31 = 8.77 \text{ dB}$$

Esto tiene sentido porque:
$$20\log_{10}\left(\frac{85}{31}\right) = 20\log_{10}(2.74) = 8.76 \text{ dB}$$ ✓

\textbf{Regla práctica:} Duplicar la frecuencia añade 6 dB de FSL.

\subsection{Pregunta 1f: Atenuación Atmosférica y Lluvia}

\subsubsection{Atenuación atmosférica en cielo claro}

La atmósfera absorbe señales RF principalmente por:

\textbf{1. Vapor de agua (H₂O):}
\begin{itemize}
    \item Picos de resonancia: 22.2 GHz, 183 GHz, 325 GHz
    \item A 31 GHz: atenuación baja ($\sim$0.05-0.1 dB/km)
    \item A 85 GHz: en ventana atmosférica, pero atenuación moderada ($\sim$0.3-0.5 dB/km)
\end{itemize}

\textbf{2. Oxígeno molecular (O₂):}
\begin{itemize}
    \item Picos de resonancia: 60 GHz (muy fuerte), 118 GHz
    \item A 31 GHz: atenuación muy baja
    \item A 85 GHz: alejado del pico de 60 GHz, moderada
\end{itemize}

\textbf{Del gráfico proporcionado (Figura 1 o anexo):}

\begin{center}
\fcolorbox{blue}{blue!10}{
\begin{minipage}{0.85\textwidth}
\textbf{Lectura del gráfico:}
\begin{itemize}
    \item \textbf{Banda Ka (31 GHz):} $\alpha_{atm} \approx 0.08$ dB/km
    \item \textbf{Banda E (85 GHz):} $\alpha_{atm} \approx 0.4$ dB/km
\end{itemize}
\textbf{Para distancia $L = 40$ km:}
\begin{itemize}
    \item $A_{atm,Ka} = 0.08 \times 40 = 3.2$ dB
    \item $A_{atm,E} = 0.4 \times 40 = 16$ dB
\end{itemize}
\end{minipage}
}
\end{center}

\textbf{Nota:} En banda E, la atenuación atmosférica es SIGNIFICATIVA incluso en cielo claro.

\subsubsection{Atenuación por lluvia}

La lluvia causa atenuación por \textbf{scattering de Mie} cuando el tamaño de las gotas es comparable a la longitud de onda:

\textbf{Longitudes de onda:}
\begin{itemize}
    \item $\lambda_{Ka} = c/f = 3 \times 10^8 / 31 \times 10^9 = 9.68$ mm
    \item $\lambda_E = 3 \times 10^8 / 85 \times 10^9 = 3.53$ mm
\end{itemize}

\textbf{Tamaño típico de gotas de lluvia:} 0.5-5 mm

\textbf{Observación:} A 85 GHz, $\lambda \approx$ tamaño de gotas $\Rightarrow$ scattering MÁXIMO!

\textbf{Fórmula empírica ITU-R P.838:}
$$\alpha_R = k \times R^{\alpha} \text{ [dB/km]}$$

Donde $R$ = tasa de precipitación [mm/h], y $k$, $\alpha$ dependen de frecuencia y polarización.

\textbf{Del gráfico (Fig. 2) para R = 25 mm/h:}

\begin{center}
\fcolorbox{red}{red!10}{
\begin{minipage}{0.85\textwidth}
\textbf{Atenuación específica por lluvia:}
\begin{itemize}
    \item \textbf{Banda Ka (31 GHz):} $\alpha_{R,Ka} \approx 3$ dB/km
    \item \textbf{Banda E (85 GHz):} $\alpha_{R,E} \approx 15$ dB/km (¡5× más!)
\end{itemize}
\end{minipage}
}
\end{center}

\textbf{¿Por qué tanta diferencia?}

$$\frac{\alpha_{R,E}}{\alpha_{R,Ka}} = \frac{15}{3} = 5$$

Esto es porque $\alpha_R \propto f^{2-3}$ aproximadamente en régimen de Mie.

\subsubsection{Longitud efectiva del trayecto en lluvia}

No toda la distancia del enlace atraviesa lluvia. Solo la parte atmosférica (troposfera).

\textbf{Geometría:}
\begin{itemize}
    \item Altura capa de lluvia: $h_R = 4$ km (típico en climas templados)
    \item Ángulo de elevación: $\varepsilon = 30°$ (calculado anteriormente)
    \item Trayecto en lluvia: $L_{eff} = \frac{h_R}{\sin(\varepsilon)}$
\end{itemize}

\textbf{CÁLCULO:}
$$L_{eff} = \frac{4 \text{ km}}{\sin(30°)} = \frac{4}{0.5} = \boxed{8 \text{ km}}$$

\textbf{Atenuación total por lluvia:}
$$A_{rain} = \alpha_R \times L_{eff}$$

\textbf{Banda Ka:}
$$A_{rain,Ka} = 3 \text{ dB/km} \times 8 \text{ km} = \boxed{24 \text{ dB}}$$

\textbf{Banda E:}
$$A_{rain,E} = 15 \text{ dB/km} \times 8 \text{ km} = \boxed{120 \text{ dB}}$$

\textbf{¡CRÍTICO!} 120 dB es una atenuación ENORME. Equivale a un factor de $10^{12}$ (un billón).

\subsection{Pregunta 1g: Potencia Recibida en Cielo Claro}

\subsubsection{Balance completo del enlace}

$$P_R = \text{EIRP} + G_{RX} - A_{FSL} - A_{atm} \quad \text{[dBm]}$$

\subsubsection{Banda Ka - Cielo claro}

\textbf{PASO 1:} Reunimos todos los valores:
\begin{align*}
\text{EIRP}_{Ka} &= 54.14 \text{ dBm}\\
G_{RX} &= 24.13 \text{ dBi} \quad \text{(antenas idénticas)}\\
A_{FSL,Ka} &= 154.31 \text{ dB}\\
A_{atm,Ka} &= 3.2 \text{ dB}
\end{align*}

\textbf{PASO 2:} Calculamos:
$$P_{R,Ka} = 54.14 + 24.13 - 154.31 - 3.2$$
$$P_{R,Ka} = 78.27 - 157.51 = \boxed{-79.24 \text{ dBm}}$$

\subsubsection{Banda E - Cielo claro}

\textbf{PASO 1:} Valores:
\begin{align*}
\text{EIRP}_E &= 64.03 \text{ dBm}\\
G_{RX} &= 33.90 \text{ dBi}\\
A_{FSL,E} &= 163.08 \text{ dB}\\
A_{atm,E} &= 16 \text{ dB}
\end{align*}

\textbf{PASO 2:} Calculamos:
$$P_{R,E} = 64.03 + 33.90 - 163.08 - 16$$
$$P_{R,E} = 97.93 - 179.08 = \boxed{-81.15 \text{ dBm}}$$

\textbf{Comparación interesante:}

A pesar de que banda E tiene:
\begin{itemize}
    \item +8.76 dB más FSL
    \item +12.8 dB más atenuación atmosférica
    \item = +21.56 dB más pérdidas totales
\end{itemize}

Solo recibe 1.91 dB menos potencia que Ka, porque tiene:
\begin{itemize}
    \item +9.89 dB más EIRP (mayor ganancia antena TX)
    \item +9.77 dB más ganancia antena RX
    \item = +19.66 dB más ganancia total
\end{itemize}

Balance: $-21.56 + 19.66 = -1.9$ dB ✓

\subsection{Pregunta 1h: Temperatura de Ruido y Potencia de Ruido}

\subsubsection{Temperatura equivalente de ruido del receptor}

El factor de ruido $F$ se relaciona con la temperatura equivalente de ruido por:

$$T_{eq,RX} = T_0(F - 1)$$

Donde $T_0 = 290$ K (temperatura de referencia estándar).

\textbf{PASO 1:} Convertir $F$ de dB a lineal:

\textbf{Banda Ka:} $F = 3$ dB
$$F_{lineal} = 10^{3/10} = 10^{0.3} = 1.995$$

$$T_{eq,Ka} = 290 \times (1.995 - 1) = 290 \times 0.995 = \boxed{288.6 \text{ K}}$$

\textbf{Banda E:} $F = 6$ dB
$$F_{lineal} = 10^{6/10} = 10^{0.6} = 3.981$$

$$T_{eq,E} = 290 \times (3.981 - 1) = 290 \times 2.981 = \boxed{864.5 \text{ K}}$$

\textbf{Observación:} Receptor de banda E es mucho más ruidoso (3× más temperatura).

\subsubsection{Temperatura de antena}

La antena capta ruido del medio ambiente. Para antena apuntando al cielo:

\begin{itemize}
    \item \textbf{Cielo claro:} $T_A \approx 30-50$ K (ruido cósmico + atmósfera)
    \item \textbf{Con lluvia:} $T_A \approx 260-280$ K (temperatura física de las gotas)
    \item \textbf{Tierra/edificios:} $T_A \approx 290$ K
\end{itemize}

\textbf{Usaremos:} $T_A = 30$ K para cielo claro

\subsubsection{Temperatura total del sistema}

$$T_{total} = T_{eq,RX} + T_A$$

\textbf{Banda Ka - Cielo claro:}
$$T_{total,Ka} = 288.6 + 30 = \boxed{318.6 \text{ K}}$$

\textbf{Banda E - Cielo claro:}
$$T_{total,E} = 864.5 + 30 = \boxed{894.5 \text{ K}}$$

\subsubsection{Potencia de ruido}

La potencia de ruido térmico en el receptor es:

$$N_R = k \cdot T_{total} \cdot B \quad \text{[W]}$$

Donde:
\begin{itemize}
    \item $k = 1.38 \times 10^{-23}$ J/K (constante de Boltzmann)
    \item $T_{total}$ = temperatura total [K]
    \item $B$ = ancho de banda [Hz]
\end{itemize}

\textbf{Banda Ka:} $B = 297$ MHz $= 297 \times 10^6$ Hz

$$N_{R,Ka} = 1.38 \times 10^{-23} \times 318.6 \times 297 \times 10^6$$
$$N_{R,Ka} = 1.306 \times 10^{-12} \text{ W} = 1.306 \times 10^{-9} \text{ mW}$$

Convertimos a dBm:
$$N_{R,Ka,dBm} = 10\log_{10}(1.306 \times 10^{-9}) = -89.84 \text{ dBm}$$

\textbf{Alternativa (fórmula rápida):}
$$N_{dBm} = -174 + 10\log_{10}(T/290) + 10\log_{10}(B_{Hz})$$
$$N_{Ka} = -174 + 10\log_{10}(318.6/290) + 10\log_{10}(297 \times 10^6)$$
$$N_{Ka} = -174 + 0.40 + 84.73 = \boxed{-88.87 \text{ dBm}}$$

\textbf{Banda E:} $B = 2000$ MHz $= 2 \times 10^9$ Hz

$$N_E = -174 + 10\log_{10}(894.5/290) + 10\log_{10}(2 \times 10^9)$$
$$N_E = -174 + 4.89 + 93.01 = \boxed{-76.10 \text{ dBm}}$$

\textbf{Nota importante:} Banda E tiene mucho más ruido por:
\begin{itemize}
    \item Mayor temperatura equivalente (+4.89 dB)
    \item Mayor ancho de banda (+8.28 dB)
    \item Total: +13.17 dB más ruido
\end{itemize}

\subsection{Pregunta 1i: Relación C/N en Cielo Claro}

\subsubsection{Cálculo de C/N}

$$\frac{C}{N} = P_R - N_R \quad \text{[dB]}$$

\textbf{Banda Ka:}
$$C/N_{Ka} = P_R - N_R = (-79.24) - (-88.87) = \boxed{9.63 \text{ dB}}$$

\textbf{Análisis del resultado:}
\begin{itemize}
    \item C/N requerido para QPSK: 13.5 dB
    \item C/N disponible: 9.63 dB
    \item \textcolor{red}{\textbf{Déficit: 3.87 dB}} ⚠️
\end{itemize}

\textbf{Conclusión Ka:} ¡El enlace NO funciona en cielo claro! Necesitaríamos:
\begin{itemize}
    \item Aumentar potencia TX: +4 dB $\Rightarrow$ 2.5 W en lugar de 1 W
    \item Antenas más grandes: +4 dB $\Rightarrow$ diámetro 0.96 m en lugar de 0.605 m
    \item Reducir distancia a $\sim$25 km
    \item Usar modulación más robusta: BPSK (requiere solo 6-7 dB)
\end{itemize}

\textbf{Banda E:}
$$C/N_E = P_R - N_R = (-81.15) - (-76.10) = \boxed{-5.05 \text{ dB}}$$

\textbf{¡NEGATIVO!} El ruido es mayor que la señal.

\textbf{Análisis del resultado:}
\begin{itemize}
    \item C/N requerido para 256-QAM: 32.5 dB
    \item C/N disponible: -5.05 dB
    \item \textcolor{red}{\textbf{Déficit: 37.55 dB}} ✗✗✗
\end{itemize}

\textbf{Conclusión E:} ¡Completamente inviable! Ni siquiera para BPSK (requiere 6-9 dB).

\textbf{Problema del enunciado:} Los valores dados parecen inconsistentes. En la práctica, enlaces a 85 GHz a 40 km requieren:
\begin{itemize}
    \item Antenas mucho más grandes (1-2 m)
    \item Potencias más altas (5-10 W)
    \item Ancho de banda menor para reducir ruido
\end{itemize}

\newpage
\section{PARTE 2: ENLACE CON LLUVIA (R = 25 mm/h)}

\subsection{Pregunta 2a: Atenuación Total con Lluvia}

\subsubsection{Componentes de atenuación}

Con lluvia, la atenuación total es:

$$A_{total} = A_{FSL} + A_{atm} + A_{rain}$$

\textbf{Banda Ka:}
\begin{align*}
A_{total,Ka} &= 154.31 + 3.2 + 24\\
&= \boxed{181.51 \text{ dB}}
\end{align*}

\textbf{Incremento respecto a cielo claro:} +24 dB

\textbf{Banda E:}
\begin{align*}
A_{total,E} &= 163.08 + 16 + 120\\
&= \boxed{299.08 \text{ dB}}
\end{align*}

\textbf{Incremento respecto a cielo claro:} +120 dB (¡factor de $10^{12}$!)

\subsection{Pregunta 2b: Potencia Recibida con Lluvia}

$$P_{R,lluvia} = \text{EIRP} + G_{RX} - A_{total}$$

\textbf{Banda Ka:}
$$P_{R,Ka} = 54.14 + 24.13 - 181.51 = \boxed{-103.24 \text{ dBm}}$$

\textbf{Banda E:}
$$P_{R,E} = 64.03 + 33.90 - 299.08 = \boxed{-201.15 \text{ dBm}}$$

\textbf{¡Impracticable!} Esta potencia es tan baja que está por debajo del ruido térmico cósmico.

\subsection{Pregunta 2c: Temperatura de Antena con Lluvia}

Cuando la antena "ve" lluvia, capta el ruido térmico de las gotas de agua:

$$T_A \approx T_{física,lluvia} \approx 280 \text{ K}$$

\textbf{Nueva temperatura total:}

\textbf{Banda Ka:}
$$T_{total,Ka} = 288.6 + 280 = \boxed{568.6 \text{ K}}$$

\textbf{Banda E:}
$$T_{total,E} = 864.5 + 280 = \boxed{1144.5 \text{ K}}$$

\subsection{Pregunta 2d: Nueva Potencia de Ruido}

\textbf{Banda Ka:}
$$N_{Ka} = -174 + 10\log_{10}(568.6/290) + 84.73$$
$$N_{Ka} = -174 + 2.92 + 84.73 = \boxed{-86.35 \text{ dBm}}$$

Incremento de ruido: $-86.35 - (-88.87) = +2.52$ dB

\textbf{Banda E:}
$$N_E = -174 + 10\log_{10}(1144.5/290) + 93.01$$
$$N_E = -174 + 5.96 + 93.01 = \boxed{-75.03 \text{ dBm}}$$

Incremento de ruido: $-75.03 - (-76.10) = +1.07$ dB

\subsection{Pregunta 2e: C/N Final con Lluvia}

\textbf{Banda Ka:}
$$C/N_{Ka} = P_R - N_R = (-103.24) - (-86.35) = \boxed{-16.89 \text{ dB}}$$

\textbf{¡NEGATIVO!} El enlace está completamente roto. La señal es 49× más débil que el ruido.

\textbf{Margen respecto a requerido:}
$$M = -16.89 - 13.5 = \textcolor{red}{\boxed{-30.39 \text{ dB}}}$$

\textbf{Banda E:}
$$C/N_E = (-201.15) - (-75.03) = \boxed{-126.12 \text{ dB}}$$

\textbf{¡CATASTRÓFICO!} La señal es $10^{12.6}$ veces más débil que el ruido.

\textbf{Margen respecto a requerido:}
$$M = -126.12 - 32.5 = \textcolor{red}{\boxed{-158.62 \text{ dB}}}$$

\newpage
\section{CONCLUSIONES Y RECOMENDACIONES}

\subsection{Resumen de Resultados}

\begin{center}
\begin{tabular}{|l|c|c|c|c|}
\hline
\textbf{Parámetro} & \textbf{Ka Claro} & \textbf{Ka Lluvia} & \textbf{E Claro} & \textbf{E Lluvia} \\
\hline
EIRP & 54.14 & 54.14 & 64.03 & 64.03 \\
$G_{RX}$ & 24.13 & 24.13 & 33.90 & 33.90 \\
$A_{FSL}$ & 154.31 & 154.31 & 163.08 & 163.08 \\
$A_{atm}$ & 3.2 & 3.2 & 16 & 16 \\
$A_{rain}$ & 0 & 24 & 0 & 120 \\
\hline
$P_R$ [dBm] & -79.24 & -103.24 & -81.15 & -201.15 \\
$N_R$ [dBm] & -88.87 & -86.35 & -76.10 & -75.03 \\
\hline
\textbf{C/N [dB]} & \textcolor{red}{9.63} & \textcolor{red}{-16.89} & \textcolor{red}{-5.05} & \textcolor{red}{-126.12} \\
C/N req & 13.5 & 13.5 & 32.5 & 32.5 \\
\textbf{Margen} & \textcolor{red}{-3.87} & \textcolor{red}{-30.39} & \textcolor{red}{-37.55} & \textcolor{red}{-158.62} \\
\hline
\textbf{Estado} & \textcolor{red}{NO} & \textcolor{red}{NO} & \textcolor{red}{NO} & \textcolor{red}{NO} \\
\hline
\end{tabular}
\end{center}

\subsection{Análisis Crítico}

\textbf{1. Banda Ka (31 GHz):}
\begin{itemize}
    \item \textcolor{orange}{Marginal en cielo claro} (déficit 3.87 dB)
    \item \textcolor{red}{Inviable con lluvia} (déficit 30.39 dB)
    \item Atenuación por lluvia moderada (24 dB)
    \item Podría funcionar con:
    \begin{itemize}
        \item Distancias más cortas (15-20 km)
        \item Antenas más grandes (0.8-1 m)
        \item Modulación adaptativa (QPSK $\leftrightarrow$ BPSK)
    \end{itemize}
\end{itemize}

\textbf{2. Banda E (85 GHz):}
\begin{itemize}
    \item \textcolor{red}{Inviable incluso en cielo claro} (C/N negativo)
    \item \textcolor{red}{Totalmente catastrófico con lluvia} (atenuación 120 dB)
    \item La atenuación por lluvia es 5× mayor que en Ka
    \item Solo viable para:
    \begin{itemize}
        \item Distancias muy cortas (<5 km)
        \item Ambientes muy secos (desiertos)
        \item Enlaces indoor o protegidos
    \end{itemize}
\end{itemize}

\subsection{Lecciones Aprendidas}

\textbf{1. Trade-off frecuencia vs propagación:}
\begin{itemize}
    \item Mayor frecuencia $\Rightarrow$ mayor ancho de banda disponible
    \item Pero también $\Rightarrow$ mayor FSL, mayor atenuación atmosférica/lluvia
    \item Antenas más pequeñas para misma ganancia
    \item Tecnología más cara y compleja
\end{itemize}

\textbf{2. Banda E (71-86 GHz):}
\begin{itemize}
    \item Enorme ancho de banda (hasta 10 GHz disponible)
    \item Pero extremadamente sensible a condiciones climáticas
    \item Solo práctica para enlaces muy cortos o indoor
    \item Aplicaciones: backhaul urbano (< 1 km), enlaces en edificios
\end{itemize}

\textbf{3. Diseño de enlaces HAPS:}
\begin{itemize}
    \item Banda Ka es límite práctico para enlaces de decenas de km
    \item Necesita diversidad en frecuencia para disponibilidad alta
    \item Modulación adaptativa esencial
    \item Control de potencia según condiciones climáticas
\end{itemize}

\vspace{1cm}
\hrule
\begin{center}
\textit{Fin de las soluciones completas y detalladas - Examen 2022 RF}\\
\textit{Autor: Soluciones paso a paso con explicaciones exhaustivas}
\end{center}

\end{document}
